\chapter{引言}

程序不再只执行固定业务规则，而从数据中估计映射 $f$。
深度学习把 $f$ 参数化为一串可微变换，用优化更新参数。

\section{日常生活中的机器学习}

日常系统把观测 $x$ 映到决策 $a$。语音识别、推荐、视觉检测都是在估计
\[
a=\hat{f}(x),\qquad \hat{f}\in\mathcal{F},
\]
其中 $\mathcal{F}$ 由数据而非手写规则确定。
$\hat{f}(x)$ 由当前输入算出：它是系统将要执行的决策或预测值（类别、分数、推荐列表中的排序分），
不是训练损失。经验增加时，同一测试分布上 $\hat{f}$ 的风险下降；固定规则程序不自动改进。

\section{机器学习中的关键组件}

学习问题写成四元组 $(\mathcal{D},\mathcal{F},L,\mathcal{A})$：
数据、模型类、目标、优化算法。四者分别提供样本、允许的映射、误差标量、以及如何改参数。

\subsection{数据}

训练集
\[
\mathcal{D}=\{(\bm{x}^{(i)},y^{(i)})\}_{i=1}^{n}
\]
提供经验。$\bm{x}^{(i)}\in\mathbb{R}^{d}$ 是第 $i$ 个样本的特征，$y^{(i)}$ 是该样本要拟合的目标
（回归时与预测同量纲的实数，分类时是类别编号）。
样本应尽可能独立同分布地来自未知 $P(\bm{x},y)$。
特征维度 $d$、标签类型和样本量 $n$ 决定 $\mathcal{F}$ 能多复杂：$n$ 是样本个数，不是误差大小。

\subsection{模型}

模型是带参数 $\bm{\theta}$ 的映射
\[
\hat{y}=f_{\bm{\theta}}(\bm{x}).
\]
线性模型对应 $\mathcal{F}=\{\bm{x}\mapsto \bm{w}^{\mathsf T}\bm{x}+b\}$：
$\bm{w}^{\mathsf T}\bm{x}+b$ 由特征加权求和再加偏置得到，回归时是与 $y$ 同量纲的点预测，不是概率。
深度模型把若干简单映射复合为 $f_L\circ\cdots\circ f_1$，最终输出仍是预测，须再经损失才变成误差。

\subsection{目标函数}

目标（损失）把预测与真值变成可比较的标量：
\[
L(\bm{\theta})
=
\frac{1}{n}\sum_{i=1}^{n}
\ell\bigl(f_{\bm{\theta}}(\bm{x}^{(i)}),y^{(i)}\bigr).
\]
$L(\bm{\theta})$ 是 $n$ 条样本损失的算术平均，表示这批预测整体有多差；它不是 $[0,1]$ 上的概率。
回归常用平方损失（均方误差，与 $y$ 的量纲平方同量纲），分类常用交叉熵（模型对真类有多惊讶，单位 nats）。
经验风险 $L(\bm{\theta})$ 是对真实风险
\[
R(\bm{\theta})=\mathbb{E}[\ell(f_{\bm{\theta}}(\bm{x}),y)]
\]
的样本近似；$R$ 是在未知分布 $P$ 上的平均损失，无法直接计算，只能用独立测试集估计。

\subsection{优化算法}

训练求解
\[
\bm{\theta}^{\ast}
\in
\operatorname*{argmin}_{\bm{\theta}} L(\bm{\theta}).
\]
$\bm{\theta}^{\ast}$ 是使训练损失最小的参数，不是预测值本身。
深度学习中 $L$ 通常非凸且无闭式解，故用梯度迭代
\[
\bm{\theta}
\leftarrow
\bm{\theta}-\eta\,\widehat{\nabla L}(\bm{\theta}),
\]
其中 $\widehat{\nabla L}$ 常由小批量估计。
第 $j$ 个分量是：该参数增加 $1$ 个单位时，损失大约增加多少；减去它乘步长 $\eta$，是为了减小 $L$。

\section{各种机器学习问题}

按是否提供标签、是否与环境交互，学习问题分成不同的条件设置。
每种设置里，模型输出与损失所代表的数并不相同。

\subsection{监督学习}

每个 $\bm{x}$ 配有 $y$，目标是估计 $P(y\mid\bm{x})$ 或 $f(\bm{x})\approx y$。
回归取 $y\in\mathbb{R}$，此时 $\hat{y}$ 是与 $y$ 同单位的点预测。
分类取 $y\in\{1,\ldots,K\}$，此时模型通常输出 $K$ 维概率，第 $k$ 个分量是“属于类 $k$”的概率，在 $[0,1]$ 且和为 $1$。

\subsection{无监督学习}

只有 $\{\bm{x}^{(i)}\}$。任务是估计结构，例如密度 $p(\bm{x})$（可以大于 $1$，本身不是概率）、
聚类划分（样本到簇编号的映射）、或低维表示
\[
\bm{z}=g(\bm{x}).
\]
$\bm{z}$ 由 $\bm{x}$ 算出，是压缩后的坐标，用来保留变化方向，不是类别概率。

\subsection{与环境互动}

数据不再一次性给定，而由策略与环境交互产生轨迹
\[
(\bm{s}_t,a_t,r_t,\bm{s}_{t+1})_{t\ge 0}.
\]
$\bm{s}_t$ 是 $t$ 时刻状态，$a_t$ 是当时采取的动作，$r_t$ 是即时标量奖励（可正可负，不是正确率），
$\bm{s}_{t+1}$ 是下一状态。分布依赖于当前策略，因此 i.i.d.\ 假设不自动成立。

\subsection{强化学习}

目标是最大化期望回报
\[
J(\pi)=\mathbb{E}_{\pi}\Bigl[\sum_{t=0}^{T}\gamma^{t} r_t\Bigr],
\quad 0\le\gamma\le 1.
\]
$J(\pi)$ 由策略 $\pi$ 下的轨迹对折扣奖励求和再取期望得到：它是“长期能拿到多少奖励”的平均水平，
与单步 $r_t$ 同量纲；$\gamma$ 接近 $1$ 时更看重远期。监督信号是延迟的标量奖励，而不是每步的正确标签。

\section{起源}

可学习机器的思想来自神经科学中的神经元加权求和、统计学中的回归与决策理论、
以及控制论中的反馈。感知机给出
\[
\hat{y}=\mathbf{1}\{\bm{w}^{\mathsf T}\bm{x}+b>0\}.
\]
先算线性打分 $\bm{w}^{\mathsf T}\bm{x}+b$（与特征同型加权后的实数），再看是否为正：
指示函数取值 $\{0,1\}$，$1$ 表示判为正类，$0$ 表示负类，不是 $[0,1]$ 上的平滑概率。
多层网络则取消线性可分限制。

\section{深度学习的发展}

关键条件是：可微复合模型、反向传播计算 $\nabla_{\bm{\theta}}L$、大规模数据与并行硬件。
$\nabla_{\bm{\theta}}L$ 的每个分量仍是“该参数增加 $1$ 时损失大约增加多少”。
深度指复合层数 $L$ 增大，$L$ 是层的个数（正整数），不是损失值；层数增大使 $\mathcal{F}$ 更丰富。

\section{深度学习的成功案例}

当 $x$ 具有局部相关结构（图像、语音、文本）且 $n$ 足够大时，
端到端训练的 $f_{\bm{\theta}}$ 在分类、检测、翻译等风险上低于手工特征加浅层模型。
这里比较的是测试风险或错误率（错误个数占总样本的比例，取值 $[0,1]$），不是训练损失本身。

\section{特点}

深度学习用端到端可微管道代替分阶段特征工程：
\[
x\;\xrightarrow{f_1}\;h_1\;\xrightarrow{f_2}\;\cdots\;\xrightarrow{f_L}\;\hat{y},
\]
$h_{\ell}=f_{\ell}(h_{\ell-1})$ 是第 $\ell$ 层算出的表示（一组激活值），$\hat{y}$ 是最终预测。
各 $f_{\ell}$ 的参数由同一 $L(\bm{\theta})$ 的梯度同时更新：每个梯度分量仍表示损失对该参数的敏感度。

\section{小结}

机器学习估计 $f_{\bm{\theta}}$，$\hat{y}=f_{\bm{\theta}}(\bm{x})$ 是预测；
深度学习把 $f_{\bm{\theta}}$ 写成深层复合，并用 $L(\bm{\theta})$ 的梯度迭代求解。
$L$ 是平均误差大小，不是概率；梯度分量是参数每增加一单位时损失的变化量。

\section{练习}

核验四元组 $(\mathcal{D},\mathcal{F},L,\mathcal{A})$ 在回归与分类中如何特化：
回归的 $\hat{y}$ 与 $y$ 同量纲，$L$ 常为均方误差；分类的 $\mathrm{softmax}$ 第 $k$ 分量是类概率，
$L$ 常为交叉熵（nats）。再核验经验风险 $L$ 是有限样本平均，真实风险 $R$ 是对未知 $P$ 的期望。
